\section{Destilación, desventaja de recursos y la oportunidad de las empresas chinas de modelos}

Al hablar de la diferencia entre los modelos chinos y estadounidenses, Yao Shunyu (姚顺宇) considera que en el último año o año y medio el gap se ha reducido, pero todavía no está claro si podrá cerrarse por completo o incluso invertirse. Menciona en particular que las empresas chinas de modelos están en desventaja en recursos de cómputo, y que esa desventaja podría forzar la aparición de estrategias interesantes, una de las cuales es la distillation.

\begin{figure}[H]
\centering
\includegraphics[width=0.88\textwidth]{frame-distillation.jpg}
\caption{Sección de destilación: se discuten la «destilación dura» y formas más inteligentes de destilación.\protect\footnotemark}
\end{figure}
\footnotetext{Intervalo de tiempo del fotograma en el video: 00:59:30--01:00:00. Este fotograma corresponde a la escena de discusión sobre destilación dura/suave.}

\begin{knowledgebox}{Lectura de la figura: cómo usa las imágenes el párrafo sobre destilación}
La imagen solo puede indicar la procedencia temporal, no los detalles del método de destilación. La información técnica sobre la destilación se debe leer en la tabla adyacente: la diferencia entre destilación dura, destilación suave, destilación de entorno y autodestilación está en la señal de entrenamiento, no en los elementos visuales.
\end{knowledgebox}


\subsection{Destilación dura y destilación suave}

Yao Shunyu divide la destilación en una «destilación dura» tosca y una destilación con mayor contenido científico. La destilación dura consiste en tomar directamente los tokens generados por un modelo grande y entrenar con ellos de manera forzada; esto plantea problemas de ética comercial y resulta intelectualmente pobre, porque muestra que el equipo no sabe en realidad qué quiere optimizar. La destilación más inteligente, en cambio, puede girar en torno a la transferencia de capacidades, el aprendizaje de preferencias, la generación de datos y la definición de tareas: no se trata de copiar las respuestas, sino de usar un modelo fuerte para ayudar a construir mejores señales de entrenamiento.

\noindent\begin{tabularx}{\textwidth}{p{0.18\textwidth}p{0.37\textwidth}X}
\toprule
Tipo & Método & Riesgo/valor \\
\midrule
Destilación dura & Recopilar directamente los tokens de salida de un modelo fuerte para entrenar & Riesgo legal y ético alto; tiende a aprender respuestas superficiales y carece de un objetivo de conducta. \\
Destilación suave & Usar un modelo fuerte para ayudar a generar tareas, retroalimentación, explicaciones o señales de preferencia & Más cercana a un problema científico: cómo transferir capacidades, procesos y evaluación a un modelo más débil. \\
Destilación de entorno & Hacer que el modelo genere trayectorias dentro del entorno de la tarea, y luego filtrarlas o calificarlas & Más valiosa para agent/coding, pero con infraestructura compleja. \\
Autodestilación & Usar el propio modelo o modelos de la misma familia para mejorar los datos de manera iterativa & Puede formar un ciclo cerrado, pero también puede amplificar errores y provocar un colapso de modo. \\
\bottomrule
\end{tabularx}

\begin{importantbox}{La esencia de la destilación no es copiar respuestas, sino diseñar señales de entrenamiento}
Si no se sabe lo que se quiere, la destilación se degrada en copiar la salida; si la conducta objetivo está clara, la destilación puede convertirse en una herramienta de ingeniería de datos y de transferencia de capacidades. La diferencia entre ambas está en si existe una definición de tarea, un criterio de evaluación y un análisis de fallos claros.
\end{importantbox}

\subsection{Por qué los robots todavía tienen una oportunidad}

En la entrevista, Yao Shunyu expresó un interés mayor por los laboratorios de robótica que por los modelos de lenguaje, porque en la robótica todavía quedan muchos problemas sin resolver bien. Los modelos de lenguaje ya no son un océano azul; si los jóvenes de verdad quieren encontrar una oportunidad, no necesariamente deben perseguir la dirección más candente, sino buscar la dirección que «nadie ha resuelto bien». La robótica, la generación multimodal, el control cuántico, entre otros, podrían ser espacios con más de océano azul.

\begin{knowledgebox}{La desventaja de recursos puede cambiar el gusto en la investigación}
Cuando el cómputo no se puede acumular sin límite, los equipos se ven obligados a dar importancia a la eficiencia, la destilación, la selección de datos, las señales de entrenamiento y la puesta en producción. La desventaja de recursos no es una ventaja en sí misma, pero puede propiciar la aparición de caminos distintos. La oportunidad de las empresas chinas de modelos podría venir precisamente de esta diferencia de caminos.
\end{knowledgebox}

\subsection{Resumen de la sección}

La destilación no es simplemente una cuestión de «robar o no», sino un problema de diseño de la señal de entrenamiento. La destilación dura implica un riesgo alto y deja al descubierto la falta de un objetivo; la destilación suave y la destilación de entorno, en cambio, pueden convertirse en una vía importante para que los equipos con recursos limitados mejoren sus capacidades.

